\chapter{Cosmología del estado genealógico global: conexión, curvatura, expansión y memoria de frontera}
\label{chap:cosmologia}

\section{La totalidad como estado compatible}

Una cosmología HMT--MD es una realización global del sistema de rutas,
fronteras y memoria. Su estado pertenece al límite de secciones compatibles;
su geometría se expresa mediante co-tétrada, conexión y métrica; su evolución
publica factor de escala, curvatura, contenido material y horizonte. La
totalidad interviene como objeto matemático porque cada sección local es una
restricción de un estado global con condiciones de cierre.

Esta idea ofrece una formulación precisa del principio de Mach. Sea
\(b:\mathcal X_\infty\to B_\infty\) la traza global de frontera y sea
\(I_v:\mathcal X_\infty\to V_v\) una respuesta inercial en un entorno
\(v\). La respuesta es machiana cuando existe
\[
 \overline I_v:\operatorname{im}(b)\to V_v
\]
tal que
\[
 I_v=\overline I_v\circ b,
\]
y cuando la restricción puramente local deja abiertas respuestas distintas.
El mapa \(\overline I_v\) existe de forma única exactamente cuando
\(I_v\) es constante sobre cada fibra de \(b\). La inercia local queda así
determinada por la clase global de cierre.

El marco de una sección se obtiene restringiendo simultáneamente el estado y
la conexión:
\[
 \mathfrak F_v(x)
 =\operatorname{Res}_v(x,\nabla_{\rm TPK}).
\]
Una geodésica afín satisface
\[
 \nabla_{\dot\gamma}\dot\gamma=0,
\]
y un marco paralelo satisface
\(\nabla_{\dot\gamma}e_a=0\). Aceleración, rotación, torsión, curvatura y
holonomía miden los defectos respecto de ese transporte inducido.

\section{\(3\) cartas de curvatura}

Los regímenes elíptico, parabólico e hiperbólico del TRIT proporcionan tres
dominios algebraicos. Su realización geométrica requiere un mapa
\[
 \mathcal M_{s,\rho}:\mathcal A_s
 \longrightarrow(\mathcal X_{s,\rho},g_{s,\rho})
\]
que preserve conexión, curvatura, orientación y datos de borde. Con esa
interfaz declarada, el atlas coordinado utiliza una carta de curvatura
negativa de tipo anti--de Sitter, una carta euclídea gnomónica y una carta de
curvatura positiva de tipo de Sitter.

La carta negativa organiza el cuadrado holográfico entre estado interior y
observables conformes de frontera. La carta euclídea realiza el descenso local
\[
 \frac{\mathcal G_L}{\mathcal I_L}=1
\]
del principio de equivalencia. La carta positiva recibe la expansión
acelerada cuando su dinámica satisface
\[
 \frac{\ddot a}{a}=H^2>0.
\]
El antecedente común conserva la razón areal--volumétrica
\[
 \frac{\mathcal L_{\rm ar}}{\mathcal L_{\rm vol}}
 =\frac{\pi_{\rm HMT}}{\varphi_{\rm HMT}^2}.
\]
La correspondencia entre signo algebraico y geometría física se demuestra en
cada carta mediante el mapa métrico, de modo que el atlas conserva el tipo de
cada transición.

La geometría de de Sitter proporciona un resultado exacto de comparación.
Para \(H>0\), sus foliaciones cerrada, plana y abierta poseen
\[
 a_{+1}(t)=H^{-1}\cosh(Ht),\qquad
 a_0(t)=a_*e^{Ht},\qquad
 a_{-1}(t)=H^{-1}\sinh(Ht),
\]
en sus dominios respectivos. Las tres satisfacen
\[
 \left(\frac{\dot a_k}{a_k}\right)^2+\frac{k}{a_k^2}=H^2,
 \qquad
 \frac{\ddot a_k}{a_k}=H^2,
\]
\[
 R^{(4)}=12H^2,\qquad
 R^{(3)}=\frac{6k}{a_k^2}.
\]
El signo \(k\) caracteriza la curvatura de la rebanada espacial; la
curvatura del espacio-tiempo completo permanece positiva y común.

\section{Memoria distribuida y expansión}

El transporte dodecafásico separa una parte par y una parte impar. La parte
par conserva el coste compartido por una ruta y su inversa; la impar conserva
la orientación. En la lectura cosmológica, el sector par puede realizarse
como respuesta media a la expansión y el sector impar como orientación de la
hoja.

Una memoria distribuida puede proceder de un generador local. El núcleo
\[
 K_q(m,n)=q^{|m-n|},\qquad 0<q<1,
\]
es positivo y su inversa es tridiagonal. Las correlaciones entre profundidades
arbitrarias son, por ello, el Green de una acción de primeros vecinos. Esta
identidad permite que la historia global conserve memoria mientras la ley de
actualización actúa localmente.

En una geometría homogénea, la respuesta par admite la forma constitutiva
\[
 p_{\rm eff}=p-3\zeta H.
\]
El coeficiente \(\zeta\) representa resistencia de la memoria simétrica a la
expansión. Su selección pertenece a la realización cosmológica y debe
descender del operador de transporte y de su carta dimensional.

\section{Completación (1000=729+243+27+1), datos de frontera y criterio de rebote de Einstein–Cartan}

La completación
\[
 1000=729+243+27+1
\]
organiza interior, estratos de frontera y sello. En una cosmología, este
patrón se realiza sobre un par \((M,\partial M)\) con co-tétrada, conexión,
materia y datos de borde. La frontera recibe carácter espacial, temporal o
nulo; la acción incorpora sus términos de borde y esquina; las condiciones de
unión fijan qué historias se prolongan.

Un rebote de Einstein--Cartan puede aparecer cuando la materia con espín
produce una contribución efectiva que crece como
\[
 s^2\propto a^{-6}.
\]
En una convención de densidad de masa, la ecuación adopta el esquema
\[
 H^2=\frac{8\pi G}{3}
 \bigl(\rho_{\rm m}-\rho_{\rm spin,m}\bigr)+\cdots.
\]
El punto de equilibrio
\(\rho_{\rm m}=\rho_{\rm spin,m}\) es candidato a un mínimo regular del
factor de escala cuando la acción de materia, el signo, las condiciones de
energía y el problema de frontera lo sostienen. HMT añade la obligación de
derivar la corriente de espín y el umbral desde el estado enriquecido.

Una transición hacia componentes cosmológicas descendientes se representa
mediante un cobordismo
\[
 \mathcal W:
 \Sigma_{\rm in}\longrightarrow
 \Sigma_{\rm prog}\sqcup\Sigma_{\rm desc}.
\]
La conservación de información se expresa mediante una carga cocíclica, una
isometría o una evolución unitaria. En el caso cuántico, las entropías, la
información mutua y las correlaciones sustituyen a una suma escalar simple.
Los términos de unión, la causalidad, la hiperbolicidad y el control de
singularidades forman parte del objeto.

\section{Sistema común de observables cosmológicos: expansión, lentes, crecimiento estructural y anisotropías de fondo}

Una realización destinada a explicar las discrepancias atribuidas a materia
oscura o energía oscura debe utilizar un conjunto común de parámetros para
curvas de rotación, lentes gravitatorias, crecimiento de estructura,
anisotropías del fondo cósmico e historia de expansión. Esta exigencia
convierte la cosmología de rutas en un sistema conjunto de observables
correlacionados sometido a contraste simultáneo.

El ciclo dodecafásico aporta además un selector armónico. Para
\[
 w_k=\frac1{12}
 \left[1+2\varepsilon\cos(3\phi_k-\delta)\right],
\]
la transformada discreta posee soporte únicamente en
\(m=0,\pm3\pmod{12}\). Una predicción celeste exige fijar antes del contraste
el mapa al cielo, el eje, la amplitud, la fase, la máscara y el estadístico.
La comparación ciega conserva la distinción entre el teorema armónico y su
realización cosmológica.

\begin{resultado}
La factorización machiana, el criterio de autoinercia, las tres foliaciones
de de Sitter, la memoria local con Green distribuido y el soporte armónico
dodecafásico son resultados exactos en sus dominios. El atlas métrico, la
ecuación efectiva de escala, el rebote y el cobordismo constituyen interfaces
físicas tipadas con condiciones explícitas y componen la cosmología de rutas.
\end{resultado}

\begin{lecturafisica}
El universo se describe como una sección global cuya frontera condiciona los
marcos locales. La expansión transforma la memoria colectiva; la curvatura y
las ondas devuelven esa transformación a las rutas; el horizonte registra la
parte accesible de la historia.
\end{lecturafisica}

\begin{frontera}
La realización quedaría refutada si la respuesta local variara dentro de una
misma fibra global, si el mapa métrico incumpliera curvatura u orientación,
si la aceleración careciera de ecuación dinámica, si el rebote produjera una
singularidad o violara las condiciones de unión, o si un único conjunto de
parámetros resultara incompatible con expansión, lentes, estructura y fondo
cósmico.
\end{frontera}
